\section{Arquitectura del Sistema}
\label{sec:arquitectura}

\subsection{Visión General}
La arquitectura diseñada se basa en microservicios dockerizados ejecutados en una Raspberry Pi. El sistema está dividido en cinco bloques funcionales principales: Entrada/Salida (Proxy), API Gateway (Backend Flask), Motor de Procesamiento (Worker), Capa de Datos (MongoDB y Redis) y Capa de Observabilidad (Tridente de Prometheus, Loki y Grafana).

\subsection{Exposición Pública y Proxy Inverso}
Para recibir los Webhooks de Telegram de forma segura sin abrir puertos en el router, se ha integrado \textbf{Ngrok}, que levanta un túnel TCP cifrado. 
El tráfico entrante es recibido por un contenedor \textbf{Nginx} (puerto 8080) configurado como Reverse Proxy.
El archivo \texttt{nginx.conf} rutea el tráfico basándose en subrutas:
\begin{itemize}
    \item \texttt{/}: Redirige al backend de Flask.
    \item \texttt{/grafana/}: Redirige al panel de Grafana (puerto 3000). Se emplea la directiva \texttt{proxy\_redirect} para interceptar re-escrituras locales generadas por el propio Grafana y adaptarlas al contexto del túnel público.
    \item \texttt{/rabbitmq/}: Redirige a la interfaz de administración de RabbitMQ (puerto 15672). Se inyecta la cabecera \texttt{Authorization: Basic} en tiempo de vuelo para realizar un auto-login seguro (Zero Trust interno).
\end{itemize}

\subsection{Desacoplamiento con RabbitMQ}
Dado que el bot puede recibir un pico repentino de mensajes o fotos pesadas, y la API de Gemini tiene tiempos de respuesta de hasta varios segundos, procesar la IA sincrónicamente colapsaría el servidor web.
Por ello, el backend en Flask (Productor) actúa únicamente como pasarela. Cuando recibe un JSON de Telegram, lo publica inmediatamente en una cola gestionada por RabbitMQ en la nube (CloudAMQP) y devuelve un código HTTP 200 a Telegram.
El script \texttt{worker.py} (Consumidor) está subscrito a dicha cola de manera asíncrona, consumiendo los mensajes de uno en uno (estrategia Prefetch Count 1) garantizando la estabilidad del sistema por mucho que aumente el tráfico.

\subsection{Arquitectura de Contenedores (Docker Compose)}
La infraestructura se levanta de manera declarativa con un \texttt{docker-compose.yml} compuesto por 9 servicios conectados en una red puente privada (\texttt{tfg\_network}):
\texttt{ngrok, nginx, backend\_telegram, worker, rabbitmq, redis, prometheus, grafana, loki, promtail, node-exporter, cadvisor}.
Todos los contenedores tienen configuradas políticas de reinicio (\texttt{restart: unless-stopped}) asegurando alta disponibilidad tras cortes de luz en la Raspberry Pi.
